% TOP_PROBLEM: {"id": 7800013, "title": "Bose-Einstein Condensation in Continuum Models", "category_id": 16, "category": {"id": 16, "name": "physics", "display_name": "Mathematical Physics", "description": "Problems at the intersection of mathematics and physics.", "slug": "physics", "order_index": 16, "created_at": "2026-07-31T15:26:25.671Z"}, "status": "open"}
% FIRST_POSTED: 2026-09-27
% ATTEMPT_STATUS: unresolved
% !TeX program = lualatex
\documentclass[11pt]{article}
\usepackage[margin=25mm]{geometry}
\usepackage{fontspec}
\setmainfont{Latin Modern Roman}
\usepackage{amsmath,amssymb,mathtools}
\usepackage{unicode-math}
\setmathfont{Latin Modern Math}
\usepackage{xurl,hyperref}
\hypersetup{colorlinks=true,urlcolor=blue}
\setcounter{secnumdepth}{0}
\setlength{\parindent}{0pt}
\setlength{\parskip}{5pt}
\setlength{\emergencystretch}{3em}
\begin{document}
\section*{\texorpdfstring{Bose-Einstein Condensation in Continuum Models}{Bose-Einstein Condensation in Continuum Models}}\label{7800013}
\subsection{Definitions and mathematical statement}
For $N$ identical bosons in a box $\Lambda\subset{\mathbb{R}}^3$, a standard interacting Hamiltonian is
\[
 H_{N,\Lambda}=\sum_{i=1}^N-\Delta_{x_i}+\sum_{i<j}v(x_i-x_j)
\]
on symmetric wave functions. The one-particle density matrix is $\gamma(x,y)=\langle a^*(x)a(y)\rangle$. Off-diagonal long-range order means that the thermodynamic-limit kernel retains a nonzero long-distance component; in translation-invariant situations this is expressed as $\lim_{|x-y|\to\infty}\gamma(x,y)>0$. The limit sends $N,|\Lambda|\to\infty$ at fixed density before long-distance analysis. The catalogue does not specify the interaction class, ensemble, boundary conditions, or positive-temperature range. These are necessary additional model data, not grounds for claiming the assertion for every interaction.
\par\smallskip\textit{Formulation and concept references:} \href{https://doi.org/10.1007/3-7643-7338-9}{[S1]}.

\subsection{Short English statement}
Can low-temperature condensation be rigorously established for a specified interacting three-dimensional continuum Bose gas in the thermodynamic limit?
\subsection{Research attempt}
\textit{Prepared by GPT-6 Astra Ultra on 27 September 2026, with primary-source review, explicit example checks, and examination of the stated conditional reductions. This is an unresolved research attempt.}

\paragraph{A specified model and source audit (27 September 2026).}
Fix units in which the kinetic energy is $-\Delta$. Use periodic cubes
$\Lambda_L=(\mathbb R/L\mathbb Z)^3$, the canonical Gibbs state at
inverse temperature $\beta$, and $N=\lfloor\rho L^3\rfloor$ with
$\rho=10^{-3}$. Fix the nonnegative soft-sphere interaction
\[
 v(x)=1_{\{|x|\le1\}},
\]
periodized on the cube, with $L>2$. These data are one concrete
continuum model for the catalogue question. The intended regime is
fixed finite $\beta$ sufficiently large, followed by $L\to\infty$;
the interaction is not rescaled with $N$.

The book [S1], available from its authors as [E1], distinguishes
thermodynamic condensation from other scaling limits. S\"ut\H{o} [S2]
claims a condensation theorem under simultaneous positivity conditions
on the potential and its Fourier transform, with additional cycle
analysis. That Fourier hypothesis does not hold here: in the convention
$\widehat v(k)=\int e^{-ik\cdot x}v(x)\,dx$,
\[
 \widehat v(k)=4\pi\frac{\sin |k|-|k|\cos|k|}{|k|^3},
 \qquad \widehat v(2\pi e_1)=-1/\pi.
\]
Its claimed theorem is therefore not invoked for this model, nor is
its complete proof independently certified here. K\"onig [S3]
explicitly studies a variational free-energy problem and leaves
off-diagonal long-range order outside its conclusions.

The proposed route is a depletion estimate: control occupations away
from zero momentum tightly enough that a positive fraction must remain
in the zero mode. An explicit kinetic-density bound can be proved.
The missing infrared estimate and the distinction between occupation
and pointwise long-range order will remain visible.

\paragraph{Occupation numbers and the order of limits.}
The canonical Gibbs state is translation invariant. Its one-particle
density matrix is diagonal in plane waves with momenta
$p\in(2\pi/L)\mathbb Z^3$. Write their occupations as $n_p\ge0$.
They satisfy
\[
 \sum_p n_p=N,\qquad
 \gamma_L(0,r)=L^{-3}\sum_p n_p e^{ip\cdot r},\qquad
 L^{-3}\langle T\rangle=L^{-3}\sum_p |p|^2n_p.
\]
A sufficient condition for ordinary zero-mode condensation is
$\liminf_{L\to\infty}n_0/L^3>0$. This condition refers to a
macroscopic eigenvalue of the density matrix, not a large occupation
at one fixed finite volume.

Define the finite momentum measures
$\nu_L=L^{-3}\sum_p n_p\delta_p$. A uniform kinetic-density bound
makes them tight, because their mass outside a ball of radius $R$
is at most $R^{-2}L^{-3}\langle T\rangle$.
Along a weakly convergent thermodynamic subsequence, their Fourier
transforms converge at every fixed separation. If its limit is
\[
 \nu=\rho_0\delta_0+(2\pi)^{-3}f(p)\,dp,
 \qquad f\in L^1(\mathbb R^3),
\]
the Riemann--Lebesgue lemma proves
$\gamma(0,r)\to\rho_0$ as $|r|\to\infty$.
The absolute continuity of the remaining momentum measure is an
additional hypothesis in this implication.

Without that regularity, a spatial average still detects the atom:
averaging $\gamma(0,r)$ over $[-R,R]^3$ multiplies the integrand at
momentum $p$ by $\prod_j\sin(Rp_j)/(Rp_j)$. This multiplier is
bounded by one and tends to zero for every $p\ne0$.
Dominated convergence therefore gives the limit $\nu(\{0\})$.
This averaged conclusion must not be substituted for the pointwise
long-distance limit in the catalogue. Also, infinitely many small
nonzero momentum modes may converge to an atom: a thermodynamic atom
alone does not automatically identify one finite-volume mode.

\paragraph{A rigorous kinetic-energy density estimate at positive temperature.}
Let $Z_N(\beta)=\operatorname{Tr}_{\mathrm{sym}}e^{-\beta H}$ and
$f_N(\beta)=\log Z_N(\beta)$. The constant normalized $N$-body
wavefunction has trial energy
\[
 E_{\mathrm{tr}}=\frac{N(N-1)}{2L^3}\widehat v(0),
 \qquad \widehat v(0)=4\pi/3.
\]
The spectral Jensen inequality gives
$Z_N(\beta)\ge e^{-\beta E_{\mathrm{tr}}}$.
Because $H\ge T$, the min--max principle bounds its partition function
above by the free canonical partition function. This trace comparison
does not require the exponential function to be operator monotone.

Convexity of $f_N$ gives
\[
 \langle H\rangle_\beta=-f_N'(\beta)
 \le\frac{2}{\beta}\bigl(f_N(\beta/2)-f_N(\beta)\bigr)
 \le 2E_{\mathrm{tr}}+\frac{2}{\beta}\log Z_N^{(0)}(\beta/2).
\]
For the ideal canonical gas, the zero mode fills whatever particles
are not excited. Discarding the constraint on their total number gives
\[
 Z_N^{(0)}(\tau)\le
        \prod_{p\ne0}(1-e^{-\tau|p|^2})^{-1}.
\]
At fixed $\tau>0$, the logarithm of this product divided by $L^3$
converges to
\[
 (2\pi)^{-3}\int_{\mathbb R^3}-\log(1-e^{-\tau|p|^2})\,dp
   =\frac{\zeta(5/2)}{(4\pi\tau)^{3/2}}.
\]
Expand the logarithm into its positive series and integrate each
Gaussian to obtain the equality. The logarithmic singularity at zero
is integrable; shell estimates justify the corresponding Riemann-sum
limit after excluding a small ball. Since the interaction is nonnegative,
$\langle T\rangle\le\langle H\rangle$. Consequently
\begin{equation}\label{bec:energy}
 \limsup_{L\to\infty}L^{-3}\langle T\rangle_\beta
 \le e_*(\beta):=\frac{4\pi}{3}\rho^2
       +\frac{2\zeta(5/2)}{\beta(2\pi\beta)^{3/2}}.
\end{equation}
This is an actual bound for the chosen interacting canonical model.
It does not by itself force a macroscopic eigenvalue: kinetic energy
assigns arbitrarily small costs to arbitrarily small nonzero momenta.

\paragraph{A precise conditional infrared criterion.}
Suppose, in addition, that for some fixed constants $A,B$ and
$\kappa>0$, uniformly in sufficiently large $L$,
\begin{equation}\label{bec:infrared}
 n_p\le\frac{A}{|p|}+\frac{B}{\beta|p|^2}
 \qquad(0<|p|\le\kappa).
\end{equation}
Split the excited density at $\kappa$. The high-momentum contribution
is at most $\kappa^{-2}L^{-3}\langle T\rangle$.
Both singular functions in \eqref{bec:infrared} are integrable at
zero in dimension three, and their lattice Riemann sums have limits
\[
 (2\pi)^{-3}\int_{|p|\le\kappa}\frac{dp}{|p|}
       =\frac{\kappa^2}{4\pi^2},\qquad
 (2\pi)^{-3}\int_{|p|\le\kappa}\frac{dp}{|p|^2}
       =\frac{\kappa}{2\pi^2}.
\]
One can control the omitted ball uniformly by grouping lattice points
in shells; its contribution is $O(\varepsilon^2)$ or
$O(\varepsilon)$ respectively. Thus
\begin{equation}\label{bec:depletion}
 \liminf_{L\to\infty}\frac{n_0}{L^3}
 \ge\rho-\frac{A\kappa^2}{4\pi^2}
          -\frac{B\kappa}{2\pi^2\beta}
          -\frac{e_*(\beta)}{\kappa^2}.
\end{equation}
Every term is measured per unit volume, with $\rho$ fixed before
the thermodynamic limit.

There is a concrete nonempty budget for this sufficient criterion.
For $A=1$, $\kappa=1/10$, and the chosen $\rho=10^{-3}$, the
temperature-independent depletion allowance is
\[
 \frac{10^{-2}}{4\pi^2}
 +\frac{(4\pi/3)10^{-6}}{10^{-2}}
 <0.000673<\rho.
\]
Therefore a proof of \eqref{bec:infrared} with these constants and
some $B$ uniform for all sufficiently large $\beta$ would imply
positive zero-mode condensation at sufficiently low positive temperature.
This is a conditional numerical budget, not a claimed infrared bound.
The coefficients in a physically motivated approximation cannot be
inserted into it as rigorous estimates without a uniform error theorem.

\paragraph{Why the easy substitutes do not provide that estimate.}
For the ideal gas, integrating
$(e^{\beta|p|^2}-1)^{-1}$ gives the critical excited density
\[
 \rho_c^{(0)}(\beta)=\frac{\zeta(3/2)}{(4\pi\beta)^{3/2}}.
\]
This benchmark explains the relevance of an integrable infrared
singularity. It does not justify using ideal-gas occupations in the
interacting canonical state. Operator monotonicity of occupation
numbers under addition of a repulsive many-body interaction has not
been proved by the partition-function comparison above.

A purely thermal estimate on \emph{all} excited modes, vanishing as
$\beta\to\infty$, would also discard quantum depletion. Already for
two particles with this nonconstant interaction, the constant product
wavefunction is not an eigenvector: applying $H$ multiplies it by
$v(x_1-x_2)$. The positive finite-volume ground state therefore cannot
be exactly that product. The temperature-independent term in
\eqref{bec:infrared} leaves room for such excitations. This observation
does not prove its coefficient $A$ or its momentum dependence.

Similarly, a bound on free energy alone is not an estimate on every
$n_p$. A one-particle source $s n_p$ would make its derivative an
occupation number, but extracting that derivative requires bounds
uniform in the source and precise enough on both sides. No such
source-dependent control follows from \eqref{bec:energy}.
Finite-volume periodic kernels also have no literal infinite-distance
limit: first one must construct a thermodynamic correlation limit.

\paragraph{Verification, gap, and outcome.}
The local calculations check the Fourier-transform sign, Gaussian
integral constants, and the strict depletion budget. The proofs use
the canonical state throughout the interacting argument. The ideal
critical-density formula is a separately labelled comparison, not an
ensemble-equivalence theorem for the interacting gas.

The first missing estimate is \eqref{bec:infrared}, or a substitute
whose integrated depletion satisfies the strict inequality in
\eqref{bec:depletion}. Establishing the full pointwise ODLRO target
also requires a thermodynamic correlation limit and control of its
noncondensed momentum part. A successful continuation could obtain
source-dependent occupation estimates, prove a uniform low-momentum
bound, or justify the necessary correlation regularity. None is
provided by taking temperature to zero before volume to infinity.

\textbf{UNRESOLVED.} An explicit kinetic-density bound and a conditional
condensation criterion are proved for a fixed interacting model.
Neither macroscopic occupation nor pointwise ODLRO is established
unconditionally. No novelty or formal proof verification is claimed.

\subsection{Sources}
\begin{itemize}
\item[S1]E. H. Lieb, R. Seiringer, J. P. Solovej, J. Yngvason, The Mathematics of the Bose Gas, Birkhäuser, 2005.
\url{https://doi.org/10.1007/3-7643-7338-9}.
\textit{Formulation source.}
\item[S2]Bose-Einstein condensation of interacting bosons: A two-step proof.
\url{https://arxiv.org/abs/2305.18959}.
\textit{Further reference; not a proof endorsement.}
\item[S3]The free energy of the interacting Bose gas: A variational description with loops and interlacements.
\url{https://www.wias-berlin.de/people/koenig/www/CondensateInterlace.pdf}.
\textit{Further reference; not a proof endorsement.}
\item[E1]Elliott H. Lieb, Robert Seiringer, Jan Philip Solovej, and Jakob Yngvason, \emph{The Mathematics of the Bose Gas and its Condensation}, author version.
\url{https://arxiv.org/abs/cond-mat/0610117}.
\textit{Author-provided version of [S1]; consulted 27 September 2026 for the distinction between limits and models. The finite-volume energy estimate and conditional occupation bound are derived in the attempt.}
\end{itemize}
\end{document}
